\chapter{Apéndice}
\label{ch:apendice}

\section{Esquema completo del DWH}
\label{ap:ddl}

Se reproducen aquí los \texttt{DDL} de \texttt{sql/02\_dwh\_schema/} tal
como están en el repositorio. Las claves subrogadas usan \texttt{serial}
y las claves de negocio se marcan \texttt{unique}.

\begin{lstlisting}[language=SQLpostgres,
  caption={\texttt{sql/02\_dwh\_schema/00\_create\_schema.sql}: creación
  del esquema \texttt{dwh}.},
  label={lst:ap-create-schema}]
-- Esquema dwh: modelo en estrella.
create schema if not exists dwh;
\end{lstlisting}

\begin{lstlisting}[language=SQLpostgres,
  caption={\texttt{sql/02\_dwh\_schema/10\_dim\_tables.sql}: dimensiones
  del DWH.},
  label={lst:ap-dim-tables}]
-- Tablas de dimension.

drop table if exists dwh.dim_date cascade;
create table dwh.dim_date (
    date_id      int           primary key,         -- yyyymmdd
    full_date    date          not null unique,
    year         smallint      not null,
    quarter      smallint      not null,
    month        smallint      not null,
    month_name   varchar(15)   not null,
    day          smallint      not null,
    day_of_week  smallint      not null,
    day_name     varchar(15)   not null,
    week_of_year smallint      not null,
    is_weekend   boolean       not null
);

drop table if exists dwh.dim_customer cascade;
create table dwh.dim_customer (
    customer_sk  serial        primary key,
    customer_id  int           not null unique,
    full_name    varchar(200),
    email        varchar(200),
    phone        varchar(50),
    signup_date  date
);

drop table if exists dwh.dim_product cascade;
create table dwh.dim_product (
    product_sk   serial        primary key,
    product_id   int           not null unique,
    name         varchar(200),
    category     varchar(100),
    brand        varchar(100),
    sku          varchar(50),
    unit_cost    numeric(12,2),
    unit_price   numeric(12,2)
);

drop table if exists dwh.dim_store cascade;
create table dwh.dim_store (
    store_sk     serial        primary key,
    store_id     int           not null unique,
    name         varchar(200),
    address      varchar(200),
    city         varchar(100),
    postal_code  varchar(10),
    district     varchar(100),
    area_type    varchar(50),
    zone_orientation varchar(50),
    opened_date  date
);

drop table if exists dwh.dim_offer cascade;
create table dwh.dim_offer (
    offer_sk         serial      primary key,
    offer_id         int         not null unique,
    name             varchar(200),
    discount_percent numeric(5,2),
    start_date       date,
    end_date         date
);

drop table if exists dwh.dim_return_reason cascade;
create table dwh.dim_return_reason (
    reason_sk    serial        primary key,
    reason_id    int           not null unique,
    reason       text,
    active       boolean
);
\end{lstlisting}

\begin{lstlisting}[language=SQLpostgres,
  caption={\texttt{sql/02\_dwh\_schema/20\_fact\_tables.sql}: hechos del
  DWH.},
  label={lst:ap-fact-tables}]
-- Tablas de hechos.

drop table if exists dwh.fact_sales cascade;
create table dwh.fact_sales (
    sale_item_id  int           primary key,
    sale_id       int           not null,
    sale_date_id  int           not null references dwh.dim_date(date_id),
    customer_sk   int           references dwh.dim_customer(customer_sk),
    product_sk    int           references dwh.dim_product(product_sk),
    store_sk      int           references dwh.dim_store(store_sk),
    offer_sk      int           references dwh.dim_offer(offer_sk),
    quantity      int           not null,
    unit_price    numeric(12,2) not null,
    unit_cost     numeric(12,2),
    subtotal      numeric(14,2) not null,
    cost          numeric(14,2),
    margin        numeric(14,2)
);
create index ix_fs_customer on dwh.fact_sales(customer_sk);
create index ix_fs_product  on dwh.fact_sales(product_sk);
create index ix_fs_store    on dwh.fact_sales(store_sk);
create index ix_fs_date     on dwh.fact_sales(sale_date_id);

drop table if exists dwh.fact_returns cascade;
create table dwh.fact_returns (
    return_id        int           primary key,
    sale_item_id     int           references dwh.fact_sales(sale_item_id),
    return_date_id   int           not null references dwh.dim_date(date_id),
    customer_sk      int           references dwh.dim_customer(customer_sk),
    product_sk       int           references dwh.dim_product(product_sk),
    store_sk         int           references dwh.dim_store(store_sk),
    reason_sk        int           references dwh.dim_return_reason(reason_sk),
    quantity         int           not null,
    returned_value   numeric(14,2)
);
create index ix_fr_customer on dwh.fact_returns(customer_sk);
create index ix_fr_product  on dwh.fact_returns(product_sk);
create index ix_fr_date     on dwh.fact_returns(return_date_id);
\end{lstlisting}

\section{Transformaciones SQL (resumen)}
\label{ap:transformaciones}

Las transformaciones \texttt{20\_load\_dim\_customer.sql} a
\texttt{50\_load\_dim\_misc.sql} se reproducen tal cual; las dos cargas
de hechos ya aparecen en los listados~\ref{lst:fact_sales_load} y los
\texttt{LEFT JOIN} de \texttt{fact\_returns} se citan a continuación.

\begin{lstlisting}[language=SQLpostgres,
  caption={\texttt{30\_load\_dim\_product.sql}: combinación de
  \texttt{product} y \texttt{central\_product}.},
  label={lst:ap-dim-product}]
-- dim_product: combina staging.product (50) con staging.central_product (49).
truncate table dwh.dim_product restart identity cascade;

insert into dwh.dim_product (product_id, name, category, brand, sku, unit_cost, unit_price)
select  p.product_id,
        coalesce(cp.name, p.name)                                         as name,
        coalesce(cat.name, p.category)                                    as category,
        coalesce(b.name,  p.manufacturer)                                 as brand,
        cp.sku,
        cp.unit_cost,
        coalesce(cp.unit_price, p.price)                                  as unit_price
from staging.product p
left join staging.central_product cp on cp.product_id = p.product_id
left join staging.category   cat on cat.category_id = cp.category_id
left join staging.brand      b   on b.brand_id      = cp.brand_id;
\end{lstlisting}

\begin{lstlisting}[language=SQLpostgres,
  caption={\texttt{70\_load\_fact\_returns.sql}: reutilización de FKs de
  \texttt{fact\_sales}.},
  label={lst:ap-fact-returns}]
truncate table dwh.fact_returns cascade;

insert into dwh.fact_returns
       (return_id, sale_item_id, return_date_id, customer_sk, product_sk,
        store_sk, reason_sk, quantity, returned_value)
select  ri.return_id,
        ri.sale_item_id,
        to_char(ri.return_date, 'YYYYMMDD')::int  as return_date_id,
        fs.customer_sk,
        fs.product_sk,
        fs.store_sk,
        drr.reason_sk,
        ri.quantity,
        ri.quantity * fs.unit_price               as returned_value
from staging.return_item ri
left join dwh.fact_sales       fs  on fs.sale_item_id = ri.sale_item_id
left join dwh.dim_return_reason drr on drr.reason_id  = ri.reason_id;
\end{lstlisting}

\section{Listados completos de Python}
\label{ap:python}

\begin{lstlisting}[language=Python,
  caption={\texttt{src/analytics/cltv.py}.},
  label={lst:ap-cltv}]
"""Calculo de CLTV por cliente.

Formula del enunciado:
    CLTV = Ingresos_t x Margen_t x Frecuencia_t x R_t

Donde, en la ventana temporal t:
    - Ingresos_t   = ingresos netos del cliente (ventas - devoluciones).
    - Margen_t     = (ingresos - coste) / ingresos  (proporcion).
    - Frecuencia_t = n. de tickets distintos del cliente.
    - R_t          = retencion = meses_activos / meses_de_la_ventana.
"""
from __future__ import annotations
import pandas as pd

def compute_cltv(df : pd.DataFrame) -> pd.DataFrame:
    """Calcula el CLTV de cada cliente con compras."""
    df["cltv"] = (
        df["net_revenue"].astype(float)
        * df["margin_rate"].astype(float)
        * df["frequency"].astype(float)
        * df["retention_rate"].astype(float)
    )
    return df.sort_values("cltv", ascending=False).reset_index(drop=True)
\end{lstlisting}

\begin{lstlisting}[language=Python,
  caption={\texttt{src/analytics/clustering.py}.},
  label={lst:ap-clustering}]
"""Segmentacion de clientes con PCA + KMeans."""
from __future__ import annotations

from dataclasses import dataclass
from pathlib import Path

import numpy as np
import pandas as pd
from sklearn.cluster import KMeans
from sklearn.decomposition import PCA
from sklearn.preprocessing import StandardScaler


@dataclass
class ClusteringResult:
    df: pd.DataFrame                  # df + columnas pc1, pc2, cluster
    pca: PCA
    scaler: StandardScaler
    kmeans: KMeans
    explained_variance: np.ndarray
    inertia_curve: dict[int, float]   # k -> inertia, para escoger k


def _elbow(X: np.ndarray, ks=range(2, 9), random_state=42) -> dict[int, float]:
    return {k: KMeans(n_clusters=k, n_init=10, random_state=random_state).fit(X).inertia_
            for k in ks}


def run_pca_kmeans(
    df: pd.DataFrame,
    feature_cols: list[str],
    n_clusters: int = 3,
    n_components: int = 2,
    random_state: int = 42,
) -> ClusteringResult:
    """Estandariza, aplica PCA(n_components) y KMeans(n_clusters)."""
    X = df[feature_cols].astype(float).fillna(0).to_numpy()
    scaler = StandardScaler().fit(X)
    Xs = scaler.transform(X)

    pca = PCA(n_components=n_components, random_state=random_state).fit(Xs)
    Xp = pca.transform(Xs)

    inertia = _elbow(Xp, random_state=random_state)
    km = KMeans(n_clusters=n_clusters, n_init=10, random_state=random_state).fit(Xp)

    out = df.copy()
    for i in range(n_components):
        out[f"pc{i+1}"] = Xp[:, i]
    out["cluster"] = km.labels_
    return ClusteringResult(out, pca, scaler, km, pca.explained_variance_ratio_, inertia)


def save_artifacts(result: ClusteringResult, out_dir: str | Path) -> None:
    import joblib
    out = Path(out_dir)
    out.mkdir(parents=True, exist_ok=True)
    joblib.dump(result.scaler, out / "scaler.joblib")
    joblib.dump(result.pca, out / "pca.joblib")
    joblib.dump(result.kmeans, out / "kmeans.joblib")
    result.df.to_parquet(out / "customer_segments.parquet", index=False)
\end{lstlisting}

\begin{lstlisting}[language=Python,
  caption={\texttt{src/db/engine.py}.},
  label={lst:ap-engine}]
"""Conexion a Postgres y helpers para ejecutar SQL."""
from __future__ import annotations

from functools import lru_cache
from pathlib import Path

import pandas as pd
from sqlalchemy import Engine, create_engine, text

from config.settings import settings


@lru_cache(maxsize=1)
def get_engine() -> Engine:
    """Engine SQLAlchemy reutilizable (singleton)."""
    return create_engine(settings.database_url, future=True, pool_pre_ping=True)


def run_sql(sql: str, **params) -> None:
    """Ejecuta una sentencia SQL (DDL o DML) con autocommit."""
    eng = get_engine()
    with eng.begin() as conn:
        conn.execute(text(sql), params or None)


def run_sql_file(path: str | Path) -> None:
    """Ejecuta un archivo .sql completo. Asume sentencias separadas por ';'.

    Compatible con bloques DO $$ ... $$ porque ejecuta el archivo entero
    en una sola transaccion usando exec_driver_sql.
    """
    raw = Path(path).read_text(encoding="utf-8")
    eng = get_engine()
    with eng.begin() as conn:
        conn.exec_driver_sql(raw)


def read_sql(sql: str, **params) -> pd.DataFrame:
    """Devuelve un DataFrame desde una consulta."""
    eng = get_engine()
    return pd.read_sql(text(sql), eng, params=params or None)
\end{lstlisting}

\begin{lstlisting}[language=Python,
  caption={\texttt{src/etl/load.py}.},
  label={lst:ap-load}]
"""Helpers de carga: ejecutar SQL en orden y resetear esquemas."""
from __future__ import annotations

from pathlib import Path

from src.db import run_sql, run_sql_file
from src.utils import get_logger

log = get_logger("etl.load")


def run_sql_directory(directory: str | Path) -> list[Path]:
    """Ejecuta los .sql de un directorio en orden alfabetico."""
    files = sorted(Path(directory).glob("*.sql"))
    for f in files:
        log.info("Ejecutando %s", f.name)
        run_sql_file(f)
    return files
\end{lstlisting}

\section{Manifiesto de tests}
\label{ap:tests}

Los listados completos de \texttt{tests/test\_cltv.py} y
\texttt{tests/test\_dwh\_integrity.py} aparecen a continuación. En el
capítulo~\ref{ch:conclusiones} ya se describe brevemente cada test.

\begin{lstlisting}[language=Python,
  caption={\texttt{tests/test\_cltv.py}.},
  label={lst:ap-test-cltv}]
"""Validan la logica de CLTV y metricas."""
import pytest

from src.db import read_sql

try:
    read_sql("select 1")
except Exception as e:  # pragma: no cover
    pytest.skip(f"DB no disponible: {e}", allow_module_level=True)

from src.analytics import compute_cltv, compute_customer_metrics


def test_cltv_columns_present():
    df = compute_cltv()
    expected = {
        "customer_id", "full_name", "gross_revenue", "net_revenue",
        "margin_rate", "frequency", "retention_rate", "cltv",
    }
    assert expected.issubset(df.columns)


def test_cltv_formula_consistency():
    df = compute_cltv().head(50).copy()
    expected = (
        df["net_revenue"].astype(float)
        * df["margin_rate"].astype(float)
        * df["frequency"].astype(float)
        * df["retention_rate"].astype(float)
    )
    assert ((df["cltv"] - expected).abs() < 1e-6).all()


def test_metrics_aov_matches_revenue_div_frequency():
    m = compute_customer_metrics().head(100).copy()
    pos = m[m["frequency"] > 0]
    expected = pos["net_revenue"].astype(float) / pos["frequency"].astype(float)
    assert ((pos["aov"].astype(float) - expected).abs() < 1e-6).all()


def test_recency_is_non_negative():
    m = compute_customer_metrics()
    assert (m["recency_days"] >= 0).all()
\end{lstlisting}

\begin{lstlisting}[language=Python,
  caption={\texttt{tests/test\_dwh\_integrity.py}.},
  label={lst:ap-test-dwh}]
"""Validan que el DWH este coherente respecto al raw."""
import pytest

from src.db import read_sql

# Si no hay BD disponible, saltamos todos los tests de integridad.
try:
    read_sql("select 1 as ok")
except Exception as e:  # pragma: no cover
    pytest.skip(f"DB no disponible: {e}", allow_module_level=True)


def test_fact_sales_matches_raw_sale_item_count():
    raw = read_sql("select count(*) as n from public.sale_item").iloc[0]["n"]
    dwh = read_sql("select count(*) as n from dwh.fact_sales").iloc[0]["n"]
    assert int(raw) == int(dwh)


def test_fact_returns_matches_raw_return_item_count():
    raw = read_sql("select count(*) as n from public.return_item").iloc[0]["n"]
    dwh = read_sql("select count(*) as n from dwh.fact_returns").iloc[0]["n"]
    assert int(raw) == int(dwh)


def test_dim_customer_unique_business_key():
    df = read_sql(
        "select customer_id, count(*) c from dwh.dim_customer group by 1 having count(*) > 1"
    )
    assert df.empty


def test_no_dangling_dim_date_in_facts():
    bad = read_sql("""
        select count(*) as n from dwh.fact_sales f
        where not exists (select 1 from dwh.dim_date d where d.date_id = f.sale_date_id)
    """).iloc[0]["n"]
    assert int(bad) == 0


def test_revenue_matches_between_staging_and_dwh():
    stg = read_sql("select sum(subtotal) as v from staging.sale_item").iloc[0]["v"]
    dwh = read_sql("select sum(subtotal) as v from dwh.fact_sales").iloc[0]["v"]
    assert abs(float(stg) - float(dwh)) < 0.01
\end{lstlisting}
